% ─── Capítulo 1: Analizador Léxico ────────────────────────────────────────────
\chapter{FRONTEND — Análisis Léxico}

\section{Fundamentos Teóricos}

El \textbf{análisis léxico} es la primera fase del proceso de compilación. En esta etapa, el compilador lee el código fuente (una cadena de caracteres) y lo divide en unidades lógicas con significado propio llamadas \textbf{tokens}. 

Teóricamente, el análisis léxico se basa en la teoría de \textbf{Lenguajes Regulares}. Cada categoría de token (como un identificador, un número o una palabra reservada) se describe mediante una \textbf{Expresión Regular (Regex)}. Para reconocer estas expresiones en el texto de entrada, el analizador construye un \textbf{Autómata Finito Determinista (DFA)}. El DFA consume la entrada carácter por carácter y transita entre estados hasta alcanzar un estado de aceptación, logrando aislar y clasificar la subcadena leída (el \textit{lexema}).

\begin{infobox}[¿Por qué es necesario el Análisis Léxico?]
La separación entre análisis léxico y sintáctico es crucial por tres motivos:
\begin{enumerate}
    \item \textbf{Eficiencia:} Aislar tareas de bajo nivel (como descartar espacios en blanco y comentarios) permite procesarlas muy rápido, reduciendo la carga del parser sintáctico.
    \item \textbf{Simplicidad:} El parser sintáctico trabaja con unidades abstractas (tokens) en lugar de lidiar con caracteres crudos.
    \item \textbf{Portabilidad:} Las reglas específicas de codificación de caracteres y saltos de línea quedan contenidas en el lexer.
\end{enumerate}
\end{infobox}

\begin{figure}[H]
    \centering
    \begin{tikzpicture}[
        shorten >=2pt, >=stealth, on grid,
        state/.style={circle, draw=tppblue, fill=tppblue!10, text=tppdark, minimum size=1.0cm, font=\small\bfseries, thick, inner sep=0pt},
        accept/.style={circle, draw=tppgreen!80!black, double, fill=tppgreen!10, text=tppdark, minimum size=1.0cm, font=\small\bfseries, thick, inner sep=0pt, double distance=1.5pt}
    ]
        % ── Estado Inicial ──────────────────────────────────────────────
        \node[state, initial, initial text=\textbf{inicio}] (q0) at (0,0) {$q_0$};

        % ── Rama 1: IDENTIFICADORES (arriba) ─────────────────────────
        \node[accept] (q1) at (4, 2.5) {$q_1$};
        \draw[->, thick, tppblue]
            (q0) edge[bend left=15] node[above left, font=\footnotesize] {\texttt{[a-zA-Z\_]}} (q1);
        \draw[->, thick, tppblue]
            (q1) edge[loop right, looseness=6] node[right, font=\footnotesize] {\texttt{[a-zA-Z0-9\_]}} (q1);

        % ── Rama 2a: ENTEROS (medio) ──────────────────────────────────
        \node[accept] (q2) at (4, 0) {$q_2$};
        \draw[->, thick, tppblue]
            (q0) edge node[above, font=\footnotesize] {\texttt{[0-9]}} (q2);
        \draw[->, thick, tppblue]
            (q2) edge[loop above, looseness=6] node[above, font=\footnotesize] {\texttt{[0-9]}} (q2);

        % ── Rama 2b: DECIMALES (q2 → punto → q3 → dígitos → q4 acepta)
        \node[state]  (q3) at (7, 0) {$q_3$};
        \node[accept] (q4) at (10, 0) {$q_4$};
        \draw[->, thick, tppblue]
            (q2) edge node[above, font=\footnotesize] {\texttt{.}} (q3);
        \draw[->, thick, tppblue]
            (q3) edge node[above, font=\footnotesize] {\texttt{[0-9]}} (q4);
        \draw[->, thick, tppblue]
            (q4) edge[loop above, looseness=6] node[above, font=\footnotesize] {\texttt{[0-9]}} (q4);

        % ── Rama 3: COMENTARIOS // (abajo) ────────────────────────────
        \node[state]  (q5) at (4,  -2.5) {$q_5$};
        \node[state]  (q6) at (7,  -2.5) {$q_6$};
        \node[accept] (q7) at (10, -2.5) {$q_7$};
        \draw[->, thick, tppred]
            (q0) edge[bend right=15] node[below left, font=\footnotesize] {\texttt{/}} (q5);
        \draw[->, thick, tppred]
            (q5) edge node[above, font=\footnotesize] {\texttt{/}} (q6);
        \draw[->, thick, tppred]
            (q6) edge[loop above, looseness=6] node[above, font=\footnotesize] {\texttt{[{\textasciicircum}\textbackslash n]}} (q6);
        \draw[->, thick, tppred]
            (q6) edge node[above, font=\footnotesize] {\texttt{\textbackslash n}} (q7);

    \end{tikzpicture}
    \caption{%
        \textbf{Autómata Finito Determinista (DFA) del Lexer de TPP.}
        Los estados con doble borde (\textcolor{tppgreen!70!black}{verde}) son \textbf{estados de aceptación}.
        Rama superior (\textcolor{tppblue}{azul}): $q_0 \to q_1$, identificador o palabra reservada.
        Rama media: $q_0 \to q_2$, entero; $q_2 \xrightarrow{.} q_3 \to q_4$, decimal.
        Rama inferior (\textcolor{tppred}{rojo}): $q_0 \to q_5 \to q_6 \to q_7$, comentario de línea.
    }
    \label{fig:dfa_full}
\end{figure}

\section{Implementación con Flex: \texttt{lexer.l}}

En el compilador TPP, el análisis léxico se ha delegado a \textbf{Flex} (\textit{Fast Lexical Analyzer Generator}). Flex lee el archivo \texttt{lexer.l}, compila las expresiones regulares internamente convirtiéndolas primero en un Autómata Finito No Determinista (NFA) mediante la construcción de Thompson, y luego lo convierte a un DFA equivalente optimizado. Este DFA se plasma en código C (el archivo \texttt{lex.yy.c}) que exporta la función \texttt{yylex()}.

\subsection{Configuración y Expresiones Regulares Base}

El archivo inicia con código C que será inyectado directamente en el analizador generado, seguido de la declaración de macros de Flex:

\begin{lstlisting}[style=cstyle, caption={Sección de configuración del lexer}]
%{
    #include "parser.tab.h"
    #include <stdio.h>
    #include <stdlib.h>
    #include <string.h>
    
    // Contador de columna para reportar errores precisamente
    int columna = 1;
    
    // Acción ejecutada ANTES de cada token: avanza el contador de columna
    #define YY_USER_ACTION columna += yyleng;
%}

%option yylineno    // Habilita el conteo automático de saltos de línea (\n)
%option noyywrap    // No encadenar múltiples archivos fuente

// Expresiones regulares base (macros)
DIGITO [0-9]
LETRA  [a-zA-Z]
\end{lstlisting}

La directiva \texttt{YY\_USER\_ACTION} es esencial para el rastreo de errores: antes de que Flex ejecute la acción de cualquier regla, esta directiva actualiza la posición actual de la columna sumando \texttt{yyleng} (la longitud del lexema detectado).

\subsection{Reglas de Producción Léxica}

Las reglas léxicas definen el patrón a buscar y el código C a ejecutar. Flex utiliza el principio de \textbf{Longitud Máxima (Maximal Munch)}: intentará leer tantos caracteres como sea posible mientras la expresión regular siga coincidiendo. Si dos reglas coinciden con la misma longitud, se prefiere la regla que fue definida primero.

\begin{lstlisting}[style=cstyle, caption={Reglas para palabras reservadas e identificadores}]
// 1. PALABRAS RESERVADAS (Deben ir primero por el principio de prioridad)
"fun"       { return FUN; }
"si"        { return SI; }
"sino"      { return SINO; }

// 2. NÚMEROS (Uso del cierre positivo + para dígitos)
{DIGITO}+                { yylval.num_entero = atoi(yytext);  
                           return NUM_ENTERO; }
{DIGITO}+"."{DIGITO}+    { yylval.num_decimal = atof(yytext); 
                           return NUM_DECIMAL; }

// 3. IDENTIFICADORES (Una letra seguida de 0 o más letras o números)
{LETRA}({LETRA}|{DIGITO})* { yylval.cadena = strdup(yytext); 
                              return ID; }
\end{lstlisting}

\begin{notebox}
El atributo \texttt{yylval} es el puente de datos entre el Lexer y el Parser (definido como un \texttt{union} en Bison). Cuando reconocemos un número o un identificador, convertimos el lexema bruto \texttt{yytext} (un arreglo de \texttt{char}) a su valor correspondiente numérico o creamos una copia en memoria usando \texttt{strdup}, y lo almacenamos en \texttt{yylval} antes de retornar el token.
\end{notebox}

\subsection{Tabla Completa de Tokens}

\begin{table}[H]
\centering
\begin{longtable}{lll}
\toprule
\textbf{Categoría} & \textbf{Patrón Regular} & \textbf{Acción (Token/Ignorar)} \\
\midrule
\endfirsthead
\multicolumn{3}{c}{\tablename\ \thetable{} -- continuación} \\
\toprule
\textbf{Categoría} & \textbf{Patrón Regular} & \textbf{Acción (Token/Ignorar)} \\
\midrule
\endhead
\midrule
\multicolumn{3}{r}{Continúa...} \\
\endfoot
\bottomrule
\endlastfoot
Palabras reservadas & \texttt{"fun", "si", "sino"...} & \texttt{FUN, SI, SINO...} \\
Literales Enteros & \texttt{[0-9]+} & \texttt{NUM\_ENTERO} \\
Literales Decimales & \texttt{[0-9]+.[0-9]+} & \texttt{NUM\_DECIMAL} \\
Cadenas de texto & \texttt{"([\^{}\textbackslash"]|\textbackslash.)*"} & \texttt{CADENA} \\
Identificadores & \texttt{[a-zA-Z][a-zA-Z0-9]*} & \texttt{ID} \\
Operadores aritm. & \texttt{+, -, *, /, \%} & \texttt{MAS, MENOS, MULT, DIV, MOD} \\
Asignación & \texttt{=} & \texttt{ASIG} \\
Operadores relac. & \texttt{==, !=, <, >, <=, >=} & \texttt{IGUAL, DIFERENTE, ...} \\
Puntero de retorno & \texttt{->} & \texttt{FLECHA} \\
Delimitadores & \texttt{(, ), \{, \}, [, ]} & \texttt{PARI, PARD, LLAVEI, LLAVED...} \\
Puntuación & \texttt{; , :} & \texttt{PUNTOCOMA, COMA, DOS\_PUNTOS} \\
\midrule
Comentarios & \texttt{//.*} & \textit{Ignorado (skip)} \\
Espacios & \texttt{[ \textbackslash t\textbackslash n\textbackslash r]+} & \textit{Ignorado (skip)} \\
\end{longtable}
\caption{Mapeo de expresiones regulares a Tokens en TPP}
\label{tab:tokens}
\end{table}

\section{Detección y Manejo de Errores Léxicos}

Una característica robusta del analizador léxico de TPP es su capacidad de recuperarse ante caracteres no reconocidos. Se utiliza la regla universal de un punto (\texttt{.}) al final del archivo \texttt{lexer.l}. Puesto que Flex da prioridad a las reglas definidas antes o a las de mayor longitud, esta regla de \textit{fallback} o comodín solo atrapa caracteres inválidos (por ejemplo \texttt{@}, \texttt{\#} o \texttt{\_}).

\begin{lstlisting}[style=cstyle, caption={Regla general de manejo de errores léxicos}]
.  { 
    fprintf(stderr, 
        "Error Léxico en la línea %d, columna %d: Carácter inválido '%s'\n", 
        yylineno, columna - (int)yyleng, yytext);
}
\end{lstlisting}

Esta técnica permite al compilador advertir al programador con gran precisión posicional sin colgarse abruptamente, permitiendo eventualmente al parser descubrir errores de cascada en caso de omisiones léxicas severas.
